\documentclass[11pt,a4paper]{article}

\usepackage[utf8]{inputenc}

\usepackage[T1]{fontenc}

\usepackage[polish]{babel}

\usepackage{lmodern}

\usepackage{geometry}

\geometry{a4paper,margin=27mm}

\usepackage{microtype}

\usepackage[hidelinks]{hyperref}

% Portfolio: MSC-AR-P05; source: prospective ARC-XR research scope.

\title{Weryfikacja poprawności sekwencji manualnej w AR na podstawie dowodów wizualnych}

\author{Mateusz Pomianek}

\date{Wersja 0.3.0, 21 listopada 2026}

\begin{document}

\maketitle

\begin{abstract}

Przedmiotem opracowania jest Weryfikacja poprawności sekwencji manualnej w AR na podstawie dowodów wizualnych. Czy prosta pamięć sekwencji zwiększa wykrywalność pominiętych kroków względem analizy pojedynczej klatki? W analizie porównawczej zestawione zostaną dwa warianty wykrywania błędów: klasyfikacja bieżącej obserwacji i model uwzględniający stan sekwencji czynności. Głównym punktem końcowym jest czułość wykrycia pominiętego kroku przy ustalonym limicie fałszywych alarmów na sekwencję. Opis ma charakter wstępny i nie przedstawia niezweryfikowanych wyników.

\end{abstract}

\section{Wprowadzenie}

Rozwój adaptacyjnych systemów XR stwarza możliwość badania interakcji człowieka z informacją przestrzenną, ale samo zastosowanie AR, VR lub modeli sztucznej inteligencji nie przesądza o przewadze poznawczej czy ergonomicznej. W rozpatrywanym zagadnieniu, Weryfikacja poprawności sekwencji manualnej w AR na podstawie dowodów wizualnych, kluczowe jest rozdzielenie właściwości systemu od zmienności zachowania użytkownika. Porównanie należy zaprojektować tak, aby mierzyć jednoznaczne zjawisko, a nie tylko ogólną ocenę technologii.

\section{Cel pracy i pytanie badawcze}

Czy prosta pamięć sekwencji zwiększa wykrywalność pominiętych kroków względem analizy pojedynczej klatki? Najpierw oceniana jest detekcja błędu w benchmarku offline; test użyteczności przestrzennych instrukcji można przeprowadzić jako osobne badanie.

\section{Zakres i odrębny wkład naukowy}
Weryfikacja pominięcia kroku procedury na podstawie dowodów wizualnych; niezależnym mechanizmem jest pamięć kolejności stanów, a nie generowanie instrukcji przez agenta. Pomiar wykrycia brakującego kroku w porównaniu z bezpamięciowym detektorem, a nie kolejny generator instrukcji AR.

\section{Hipotezy badawcze}
\paragraph{H1.} Wariant z pamięcią kolejności wykrywa większy odsetek prawdziwych pominięć niż klasyfikator pojedynczej obserwacji przy tym samym limicie fałszywych alarmów.
\paragraph{H2.} Pamięć procedury zmniejsza mylenie kroków wykonanych w niewłaściwej kolejności z krokami ukończonymi.
\paragraph{H3.} Różnica skuteczności utrzymuje się dla uprzednio niewidzianych sekwencji i obiektów przy stałym budżecie obliczeniowym.

\section{Projekt badania i plan eksperymentów}
\subsection{A. Operacjonalizacja i kwalifikacja danych}
Weryfikacja pominięcia kroku procedury na podstawie dowodów wizualnych; niezależnym mechanizmem jest pamięć kolejności stanów, a nie generowanie instrukcji przez agenta. Jednostka wnioskowania: cała procedura i jej jawnie oznaczone kroki; kadry jednej procedury są zależne. Punkt końcowy i sposób oceny będą zamrożone przed próbą potwierdzającą. Należy udokumentować pochodzenie obserwacji, czas akwizycji, identyfikatory sesji, harmonogram oraz legalność użycia materiałów.

\subsection{B. Eksperyment porównawczy}
Niezależnie oznaczyć poprawne, niepełne i przestawione procedury; ten sam detektor obiektów jest wejściem obu metod. Estymand pierwszorzędowy: czułość wykrycia pominiętych kroków przy ustalonej liczbie fałszywych alarmów na kompletną procedurę. Przypisanie warunków, mianowniki i kryteria wyłączeń nie mogą być dobierane po przejrzeniu wyników.

\subsection{C. Kontrole i analiza odporności}
Dopasować próg fałszywych alarmów do całych procedur, zabezpieczyć etykiety przed przeciekiem z przyszłych klatek. Zasłonięte dłonie i obiekty, nagrania z przerwami, trudniejsze kroki, fałszywe ukończenia. Każdą odmowę systemu, niekompletną próbę i brak sygnału należy raportować osobno, a nie po cichu usuwać z próby.

\subsection{D. Walidacja na niezależnym materiale}
Ocenę należy wykonać na uprzednio odłożonych scenach, sesjach lub uczestnikach. Procedurowo-klastrowy bootstrap różnic czułości i alarmów; próg dobrany na niezależnym walidacyjnym podzbiorze. Model i próg zatrzymać przed analizą końcową; rejestrować również brak przewagi i działania niepożądane.

\section{Estymand, analiza statystyczna i decyzje}
Główny estymand: czułość wykrycia pominiętych kroków przy ustalonej liczbie fałszywych alarmów na kompletną procedurę. Jednostką powtarzalności jest cała procedura i jej jawnie oznaczone kroki; kadry jednej procedury są zależne. Procedurowo-klastrowy bootstrap różnic czułości i alarmów; próg dobrany na niezależnym walidacyjnym podzbiorze. H1 jest kandydatem na hipotezę potwierdzającą. H2 i H3 pozostają eksploracyjne, dopóki protokół nie określi ich hierarchii i kontroli wielokrotnych porównań. Minimalny efekt istotny praktycznie, liczebność prób, margines nie gorszości i progi jakości wyznaczyć na podstawie pilotażu i analizy precyzji, nie zgadywać wartości.

\section{Stan badań i granica nowości}
\begin{itemize}
\item AR assistance for precise manual procedures is established; independent sequential error detection requires a different endpoint. \cite{xrprior1}.
\item AI-generated AR work instructions are prior art; measure the specific temporal error/omission mechanism. \cite{xrprior2}.
\end{itemize}
Pomiar wykrycia brakującego kroku w porównaniu z bezpamięciowym detektorem, a nie kolejny generator instrukcji AR. Samo zastosowanie VLM, EEG, okulometrii lub adaptacyjnego HUD nie stanowi wystarczającej nowości międzynarodowej.

\section{Warunki bezpieczeństwa, sprzętu i zgodności}
Opis stanowi prospektywny plan, a nie prerejestrację ani zatwierdzoną procedurę. Przed zbieraniem danych od ludzi niezbędna jest właściwa weryfikacja etyczna i zgoda. Standardowy Quest 3 nie ma wbudowanego eye trackingu, a Tobii Eye Tracker 5 działa na stanowisku ekranowym. Odstępstwa, obciążenia i ograniczenia czujników muszą być ujawnione. Nie publikować surowych danych biometrycznych ani nagrań uczestników.

\section{Materiał i metody}

W analizie porównawczej zestawione zostaną dwa warianty wykrywania błędów: klasyfikacja bieżącej obserwacji i model uwzględniający stan sekwencji czynności. Wymagane jest zdefiniowanie jednostki analizy, zasad doboru przypadków, wersji scenariusza i źródeł pomiaru. Rozkład warunków oraz reguły wykluczenia prób należy określić przed pomiarem. W badaniach z udziałem ludzi niezbędna jest uprzednia właściwa ocena etyczna i świadoma zgoda. W testach wyłącznie symulacyjnych należy oddzielić prawdę symulacyjną od obserwacji empirycznej.

\subsection{Wynik główny i kryteria jakości}

Czułość wykrycia pominiętego kroku przy ustalonym limicie fałszywych alarmów na sekwencję. Pomiar powinien zachowywać jawne jednostki, znaczniki czasu i reguły oceny poprawności. Przy modelach predykcyjnych stosować rozłączne partycje oceny; przy eksperymentach interwencyjnych oddzielić skuteczność predykcji od przyczynowego wpływu interfejsu.

\section{Plan analizy}

Analiza zdarzeniowa dla niezależnych procedur i sesji, z przedziałami ufności dla czułości, alarmów i opóźnienia. W wynikach należy raportować wielkość efektu lub błąd estymacji z przedziałem niepewności, liczbą obserwacji i odsetkiem odrzuconych prób. Analizy dodatkowe wymagają jasnego oznaczenia jako eksploracyjne. Nie wybierać progów decyzyjnych ani modelu na podstawie obserwacji ze zbioru testowego.

\section{Wyniki}

W obecnej wersji nie raportuje się wyników pomiarów, tabel statystycznych ani parametrów wyznaczonych z nieudokumentowanych danych. Nie stwierdza się w tym opracowaniu, że eksperyment został wykonany lub zgłoszony do czasopisma.

\section{Dyskusja}

Interpretacja przyszłych wyników musi uwzględnić alternatywne mechanizmy i zależność wniosków od jakości danych, warunków zadania oraz ograniczeń sprzętu. Wartość potencjalnego rozwiązania zależy nie tylko od trafności modelu, lecz także od jego opóźnienia, stabilności i wpływu na zachowanie człowieka. Wyniki testów syntetycznych nie mogą być automatycznie uogólniane na populację uczestników.

\section{Ograniczenia badania}

Błąd anotacji granic kroków i częściowa widoczność dłoni mogą powodować pozorną poprawę wykrywania. Weryfikacja etapów procedury, nie grounding obiektów P04 ani faza podania wskazówki MSC-AR-P25. Żadna część szkicu nie dowodzi przeprowadzenia zbierania danych, uzyskania zgód etycznych ani gotowości manuskryptu do zgłoszenia.

\section{Wnioski}

Przedstawiono problem, kontrast i sposób jego planowanej oceny. Rozstrzygnięcie postawionego pytania wymaga przygotowania protokołu, zweryfikowanych źródeł oraz poprawnie udokumentowanych wyników. Bez tych materiałów nie można sformułować konkluzji empirycznych.

\begin{thebibliography}{9}
\bibitem{xrprior1} Evaluation of augmented reality instructions based on initial and dynamic assembly tolerance allocation schemes in precise manual assembly (2024), Advanced Engineering Informatics. \url{https://doi.org/10.1016/j.aei.2024.102954}.
\bibitem{xrprior2} Deep learning-based augmented reality work instruction assistance system for complex manual assembly (2024), CIRP Annals. \url{https://www.sciencedirect.com/science/article/pii/S0278612524000323}.
\end{thebibliography}
\end{document}

